% Options for packages loaded elsewhere
\PassOptionsToPackage{unicode}{hyperref}
\PassOptionsToPackage{hyphens}{url}
\PassOptionsToPackage{dvipsnames,svgnames,x11names}{xcolor}
%
\documentclass[
  11pt,
  letterpaper,
]{article}
\usepackage{amsmath,amssymb}
\usepackage{iftex}
\ifPDFTeX
  \usepackage[T1]{fontenc}
  \usepackage[utf8]{inputenc}
  \usepackage{textcomp} % provide euro and other symbols
\else % if luatex or xetex
  \usepackage{unicode-math} % this also loads fontspec
  \defaultfontfeatures{Scale=MatchLowercase}
  \defaultfontfeatures[\rmfamily]{Ligatures=TeX,Scale=1}
\fi
\usepackage{lmodern}
\ifPDFTeX\else
  % xetex/luatex font selection
\fi
% Use upquote if available, for straight quotes in verbatim environments
\IfFileExists{upquote.sty}{\usepackage{upquote}}{}
\IfFileExists{microtype.sty}{% use microtype if available
  \usepackage[]{microtype}
  \UseMicrotypeSet[protrusion]{basicmath} % disable protrusion for tt fonts
}{}
\makeatletter
\@ifundefined{KOMAClassName}{% if non-KOMA class
  \IfFileExists{parskip.sty}{%
    \usepackage{parskip}
  }{% else
    \setlength{\parindent}{0pt}
    \setlength{\parskip}{6pt plus 2pt minus 1pt}}
}{% if KOMA class
  \KOMAoptions{parskip=half}}
\makeatother
\usepackage{xcolor}
\usepackage[margin=0.9in]{geometry}
\setlength{\emergencystretch}{3em} % prevent overfull lines
\providecommand{\tightlist}{%
  \setlength{\itemsep}{0pt}\setlength{\parskip}{0pt}}
\setcounter{secnumdepth}{-\maxdimen} % remove section numbering
\usepackage{amsmath,amssymb}
\usepackage{microtype}
\usepackage{xurl}
\usepackage{fancyhdr}
\pagestyle{fancy}
\fancyhf{}
\fancyhead[L]{\small Discotope exposed-point irreducibility}
\fancyhead[R]{\small AI-reviewed; unrefereed}
\fancyfoot[C]{\thepage}
\renewcommand{\headrulewidth}{0pt}
\renewcommand{\footrulewidth}{0pt}
\setlength{\headheight}{14pt}
\setlength{\emergencystretch}{3em}
\pdfmapfile{+lm.map}
\pdfmapfile{+cm.map}
\pdfmapfile{+symbols.map}
\pdfmapfile{+cmextra.map}
\pdfmapfile{+latxfont.map}
\ifLuaTeX
  \usepackage{selnolig}  % disable illegal ligatures
\fi
\usepackage{bookmark}
\IfFileExists{xurl.sty}{\usepackage{xurl}}{} % add URL line breaks if available
\urlstyle{same}
\hypersetup{
  pdftitle={Irreducibility of the exposed point variety of a discotope},
  colorlinks=true,
  linkcolor={black},
  filecolor={Maroon},
  citecolor={Blue},
  urlcolor={black},
  pdfcreator={LaTeX via pandoc}}

\title{Irreducibility of the exposed point variety of a discotope}
\author{}
\date{30 September 2026}

\begin{document}
\maketitle

\textbf{Problem:} 30005473 / OWR-12697711-006\\
\textbf{Date:} 30 September 2026\\
\textbf{Status:} complete proof with an independent AI audit;
unrefereed; historical priority unconfirmed

\subsection{1. Exact target and result}\label{exact-target-and-result}

A disc is a linear image of a closed Euclidean unit ball. A discotope is
a finite Minkowski sum of such discs. For a compact convex set
\(K\subset\mathbb R^d\), write \(\operatorname{Exp}(K)\) for the points
which uniquely maximize some nonzero real linear functional, and put \[
E(K)=\overline{\operatorname{Exp}(K)}^{\mathrm{Zar},\mathbb C}\subseteq\mathbb C^d.
\] The exact target in the 2023 Oberwolfach report is irreducibility of
\(E(D)\) when the summand discs have dimension at least two and are
generic {[}1, p.~830, Conjecture 1{]}.

\textbf{Theorem 1.} Let \(D=D_1+\cdots+D_N\subseteq\mathbb R^d\), where
\(N\geq1\) and each \(D_i\) is a disc of dimension at least two. Then
\(E(D)\) is irreducible over \(\mathbb C\). No genericity or
full-dimensionality assumption is needed.

Thus Theorem 1 gives an affirmative answer to the entire stated target.
An independent AI audit found no mathematical gap; human peer review has
not been obtained. Section 5 also proves the corresponding generic
assertion for the larger variety appearing in the 2022 paper. Neither a
degree formula nor irreducibility of a complex critical locus is
asserted.

\subsection{2. Two elementary lemmas}\label{two-elementary-lemmas}

\textbf{Lemma 2 (connected domain of regular normals).} If
\(L_1,\ldots,L_N\) are real linear subspaces of \(\mathbb R^d\), each of
codimension at least two, then \[
U=\mathbb R^d\setminus\bigcup_i L_i
\] is a nonempty, open, path-connected set.

\textbf{Proof.} A finite union of proper real linear subspaces cannot
fill \(\mathbb R^d\): choose a nonzero linear form vanishing on each
subspace and take their product, which is a nonzero real polynomial.
Openness follows because the union is closed.

Take \(x,y\in U\). Each subspace \(L_i+\mathbb Rx\), and likewise
\(L_i+\mathbb Ry\), has dimension at most \(d-1\). Choose \(z\) outside
their finite union. The segment from \(x\) to \(z\) avoids every
\(L_i\). Indeed, if \((1-t)x+tz\in L_i\) for some \(0<t\leq1\), then
\(z\in L_i+\mathbb Rx\), a contradiction; its endpoint \(t=0\) is
already outside. The segment from \(y\) to \(z\) is handled in the same
way. These two segments give a path in \(U\). \(\square\)

\textbf{Lemma 3 (analytic images have irreducible Zariski closure).}
Suppose \(U\subseteq\mathbb R^d\) is nonempty, open, and connected, and
\(F:U\to\mathbb R^m\) is real analytic. The complex Zariski closure of
\(F(U)\) is irreducible.

\textbf{Proof.} Let \[
I=\{P\in\mathbb C[X_1,\ldots,X_m]:P(F(u))=0\ \text{for every }u\in U\}.
\] This ideal is proper, since \(1\notin I\). If \(PQ\in I\) and
\(P\notin I\), choose \(u_0\in U\) with \(P(F(u_0))\ne0\). By
continuity, \(P\circ F\) is nonzero on a neighborhood of \(u_0\), so
\(Q\circ F\) vanishes there. It is a complex-valued real-analytic
function on connected \(U\), and the real-analytic identity theorem
gives \(Q\circ F=0\) throughout \(U\). Hence \(Q\in I\), and \(I\) is
prime.

For completeness, the identity theorem used here is elementary. For a
real-analytic function vanishing on a nonempty open set, the set of
points where it vanishes on a neighborhood is open and is also
relatively closed: at a limit point, every derivative vanishes by
continuity, and its convergent Taylor series is therefore zero nearby.
Connectedness finishes the argument. Apply this to the real and
imaginary parts.

By the definition of Zariski closure, \(I\) is exactly the vanishing
ideal of the closure of \(F(U)\). A complex affine algebraic set with
prime vanishing ideal is irreducible. \(\square\)

The analyticity assumption is essential to this argument. Connectedness
of an arbitrary image alone would not imply irreducibility.

\subsection{3. Parametrizing exactly the exposed
points}\label{parametrizing-exactly-the-exposed-points}

Choose injective linear maps \(A_i:\mathbb R^{m_i}\to\mathbb R^d\), with
\(m_i=\dim D_i\geq2\), such that \[
D_i=A_i(B^{m_i}),\qquad Q_i=A_iA_i^{\mathsf T}.
\] Every disc of dimension \(m_i\) admits this form, by restricting a
singular-value decomposition to its nonzero singular directions. The
symmetric matrix \(Q_i\) is positive semidefinite of rank \(m_i\).

For \(u\in\mathbb R^d\), the support function of \(D_i\) is \[
h_i(u)=\max_{\|v\|\leq1}\langle A_i^{\mathsf T}u,v\rangle
      =\|A_i^{\mathsf T}u\|
      =\sqrt{u^{\mathsf T}Q_i u}.
\] If \(A_i^{\mathsf T}u\ne0\), equality in Cauchy--Schwarz gives the
unique maximizing point \[
p_i(u)=\frac{A_iA_i^{\mathsf T}u}{\|A_i^{\mathsf T}u\|}
      =\frac{Q_i u}{\sqrt{u^{\mathsf T}Q_i u}}. \tag{1}
\] If \(A_i^{\mathsf T}u=0\), every point of \(D_i\) maximizes the
functional, so the exposed face is the whole positive-dimensional disc.

The face of a Minkowski sum exposed by \(u\) is the Minkowski sum of the
faces exposed by \(u\) in the summands. To check this, every summand
satisfies \(\langle u,x_i\rangle\leq h_i(u)\); equality in their sum is
equivalent to equality term by term. Such a sum of nonempty faces is a
singleton if and only if each face is a singleton: a nonsingleton face
remains nonsingleton after fixing a point in every other face.

It follows that the normals exposing a point of \(D\) are precisely \[
U=\mathbb R^d\setminus\bigcup_{i=1}^N\ker A_i^{\mathsf T}
 =\mathbb R^d\setminus\bigcup_{i=1}^N\ker Q_i.
\] In particular, \(0\notin U\). Moreover, there is an exact equality of
sets \[
\operatorname{Exp}(D)=F(U),\qquad
F(u)=\sum_{i=1}^N\frac{Q_i u}{\sqrt{u^{\mathsf T}Q_i u}}, \tag{2}
\] where every square root denotes the positive real square root. This
is not merely a parametrization of a dense subset: every exposed point
is covered, because its exposing normal must give a singleton face in
each summand.

\subsection{4. Proof of Theorem 1}\label{proof-of-theorem-1}

Each \(\ker A_i^{\mathsf T}\) has codimension \(m_i\geq2\). Lemma 2
shows that the domain \(U\) in (2) is nonempty, open, and connected. On
\(U\), every quadratic form \(u^{\mathsf T}Q_i u\) is strictly positive.
Since \(t\mapsto t^{-1/2}\) is real analytic on \((0,\infty)\), the map
\(F\) is real analytic on all of \(U\).

By (2) and Lemma 3, the complex Zariski closure of the exposed points is
irreducible. \(\square\)

This proof does not require the quadratic radicals to be independent,
the parametrization to be injective, its image to be smooth, or its
differential to have any specified rank. Repeated discs, coincident
spans, lower-dimensional sums, and normals exposing the same point are
all allowed.

\subsection{5. Relation to the purely nonlinear
part}\label{relation-to-the-purely-nonlinear-part}

There is a substantive distinction between the sources. The 2023 report
formulates its conjecture for \(E(D)\) {[}1{]}. The 2022 paper defines
\[
D^{\partial}=\left\{x_1+\cdots+x_N:x_i\in\partial_{L_i}D_i\right\},
\quad L_i=\operatorname{span}D_i,
\quad S(D)=\overline{D^{\partial}\cap\partial D}^{\mathrm{Zar},\mathbb C},
\] where each summand boundary is relative to its span. Its Conjecture
8.2 concerns \(S(D)\), potentially a larger algebraic set {[}2,
Definition 3.3 and Conjecture 8.2{]}. The following separate argument
bridges that distinction in the generic case.

\textbf{Theorem 4.} Assume \(D\) is full-dimensional in \(\mathbb R^d\),
every summand has dimension at least two, and the spans satisfy \[
\dim\left(\sum_{i\in J}L_i\right)
 =\min\left(d,\sum_{i\in J}\dim L_i\right)
\quad\text{for every }J\subseteq\{1,\ldots,N\}. \tag{GP}
\] Then \(S(D)=E(D)\). In particular \(S(D)\) is irreducible.

\textbf{Proof.} Every exposed point has the form (2), with each
\(p_i(u)\in\partial_{L_i}D_i\), so \(E(D)\subseteq S(D)\).

Conversely, let \(x=\sum_i x_i\in D^{\partial}\cap\partial D\), with
\(x_i\in\partial_{L_i}D_i\). A supporting hyperplane at \(x\) supplies a
nonzero normal \(u\) with \(\langle u,x\rangle=h_D(u)\). Since support
functions add, each \(x_i\) must maximize \(u\) over \(D_i\).

Let \(J=\{i:A_i^{\mathsf T}u=0\}\). For \(i\notin J\), uniqueness gives
\(x_i=p_i(u)\). The sum \(\sum_{i\in J}L_i\) lies in the proper
hyperplane \(u^\perp\). By (GP), \[
\sum_{i\in J}\dim L_i=\dim\left(\sum_{i\in J}L_i\right)\leq d-1,
\] so those subspaces form a direct sum.

Write \(x_j=A_jv_j\) for \(j\in J\), with \(\|v_j\|=1\). The direct-sum
property says that the concatenated column matrix \(A_J=[A_j]_{j\in J}\)
is injective. Its transpose is consequently surjective. Choose
\(w\in\mathbb R^d\) such that \[
A_j^{\mathsf T}w=v_j\quad\text{for every }j\in J. \tag{3}
\] If \(J\) is empty, the point \(x\) is already exposed; otherwise take
\(u_\varepsilon=u+\varepsilon w\) with \(\varepsilon>0\). Equation (3)
gives \[
A_j^{\mathsf T}u_\varepsilon=\varepsilon v_j,
\qquad p_j(u_\varepsilon)=A_jv_j=x_j\quad(j\in J).
\] For \(i\notin J\), \(A_i^{\mathsf T}u_\varepsilon\) stays nonzero
when \(\varepsilon\) is sufficiently small, and
\(p_i(u_\varepsilon)\to p_i(u)=x_i\). Thus \(u_\varepsilon\in U\) and \[
F(u_\varepsilon)\longrightarrow x.
\] Every point of \(D^{\partial}\cap\partial D\) is therefore a
Euclidean limit of exposed points. A complex algebraic set is Euclidean
closed, so these points lie in \(E(D)\). Taking Zariski closures gives
\(S(D)\subseteq E(D)\), and Theorem 1 finishes the proof. \(\square\)

Condition (GP) holds on a nonempty Zariski-open set of choices of
full-rank matrices \(A_i\) of prescribed sizes. For each subset \(J\),
maximal possible rank of its concatenated matrix is a nonempty
Zariski-open condition. The parameter space of the matrices is
irreducible, so the finite intersection of these nonempty open sets is
nonempty. This is the genericity convention of {[}2, Section 2{]}. For
types which span a smaller ambient space generically, one may instead
work in their linear span; the published convention assumes a
full-dimensional discotope.

\subsection{6. Boundary checks and limits of the
claim}\label{boundary-checks-and-limits-of-the-claim}

\begin{enumerate}
\def\labelenumi{\arabic{enumi}.}
\tightlist
\item
  \textbf{Single disc.} For \(N=1\), the conclusion recovers
  irreducibility of its complexified ellipsoidal boundary for dimension
  at least two, including discs in a proper subspace.
\item
  \textbf{Repeated discs.} Taking several copies of the same disc leaves
  the normal domain connected. No algebraic independence of radicals is
  needed.
\item
  \textbf{One-dimensional summands cannot be admitted indiscriminately.}
  The interval \([-1,1]\subset\mathbb R\) has exposed-point variety
  \(\{-1,1\}\), which is reducible. More geometrically, the sum in
  \(\mathbb R^2\) of the unit disc and the segment \([-a,a]e_1\),
  \(a>0\), has exposed points on two circular caps. Their complex
  Zariski closure is the union of the distinct circles \((X-a)^2+Y^2=1\)
  and \((X+a)^2+Y^2=1\). The removed kernel is now a hyperplane and the
  analytic domain is disconnected.
\item
  \textbf{The varieties \(S(D)\) and \(E(D)\) really can differ outside
  generic position.} Embed two identical unit discs in the \(xy\)-plane
  in \(\mathbb R^3\), and add the unit disc in the \(xz\)-plane. Let
  \(u=e_3\). The point \(e_3\) is a sum of boundary points
  \(e_1+(-e_1)+e_3\), and lies in the exposed face \(2B^2_{xy}+e_3\),
  hence in \(D^{\partial}\cap\partial D\). But \(E(D)\) is contained in
  the zero set of \[
  P(X,Y,Z)=(X^2+Y^2+Z^2-5)^2-4(4-Y^2)(1-Z^2).
  \] For this example, (2) gives \(X=2u_1/r+u_1/s\), \(Y=2u_2/r\),
  \(Z=u_3/s\), where \(r^2=u_1^2+u_2^2\) and \(s^2=u_1^2+u_3^2\). Direct
  substitution verifies \(P\circ F=0\), while \(P(0,0,1)=16\). This
  polynomial identity is checked independently by the accompanying
  exact-arithmetic script. Consequently \(e_3\notin E(D)\). Theorem 4
  uses genericity for a genuine reason.
\item
  \textbf{Not proved here.} No conclusion is drawn about the full
  complex critical locus of the addition map, birationality, or the
  degrees sought in the source papers. These are different questions.
\end{enumerate}

\subsection{References}\label{references}

{[}1{]} Chiara Meroni, \emph{Two convex conjectures for different
flavours}, in \emph{New Directions in Real Algebraic Geometry},
Oberwolfach Reports 15/2023, pp.~829--832; especially p.~830, Definition
1 and Conjecture 1. \href{https://doi.org/10.4171/OWR/2023/15}{DOI
10.4171/OWR/2023/15}.
\href{https://publications.mfo.de/bitstream/handle/mfo/4031/OWR_2023_15.pdf?sequence=4}{Official
report PDF}.

{[}2{]} Fulvio Gesmundo and Chiara Meroni, \emph{The Geometry of
Discotopes}, Le Matematiche 77(1) (2022), 143--171.
\href{https://doi.org/10.4418/2022.77.1.8}{DOI 10.4418/2022.77.1.8}.
Sections 2--3 and Conjecture 8.2.
\href{https://arxiv.org/abs/2111.01241}{arXiv:2111.01241};
\href{https://lematematiche.dmi.unict.it/index.php/lematematiche/article/download/2338/1156/6734}{publisher
PDF}.

\subsection{Verification status}\label{verification-status}

The proof is self-contained after elementary real-analytic uniqueness
and the prime-ideal criterion for irreducibility, both made explicit
above. The accompanying computations are sanity checks, not the
justification for the general result. An independent AI mathematical
audit passed this proof on 30 September 2026; the reviewed snapshot and
full report are preserved in the source bundle. This is not a formal
proof certificate or human peer review, and historical priority remains
unconfirmed. Only status wording and reference formatting were updated
after the audit; the mathematical content is unchanged.

\end{document}
